\begin{definition}
Let $\Omega\subseteq\mathbb{C}$ and let $f,a,b\colon\mathbb{C}\to\mathbb{C}$.
Identify $\mathbb{C}$ with $\mathbb{R}^2$ by $z=x+iy$ and write
$\partial_0=\partial/\partial x$, $\partial_1=\partial/\partial y$.  We say
that $a$ and $b$ are the \emph{weak Wirtinger derivatives of $f$ on $\Omega$}
if there are functions $f_x,f_y\colon\mathbb{C}\to\mathbb{C}$ such that
\begin{enumerate}
\item[(a)] $f$, $f_x$ and $f_y$ belong to $L^1_{\mathrm{loc}}(\Omega)$ in the
sense of \noderef{MemLpLocOn};
\item[(b)] $f_x$ is the distributional $\partial_0$-derivative of $f$ on
$\Omega$ and $f_y$ is the distributional $\partial_1$-derivative of $f$ on
$\Omega$, both in the sense of \noderef{IsWeakPartialOn};
\item[(c)] for every $z\in\mathbb{C}$,
\[
  a(z)=\tfrac12\bigl(f_x(z)-i\,f_y(z)\bigr),
  \qquad
  b(z)=\tfrac12\bigl(f_x(z)+i\,f_y(z)\bigr).
\]
\end{enumerate}
We then write $a=\partial f/\partial z$ and $b=\partial f/\partial\bar z$ in
the weak sense, and also $a=f_z$, $b=f_{\bar z}$.

This is the pair of operators
\[
  \frac{\partial}{\partial z}=\frac12\Bigl(\frac{\partial}{\partial x}-i\frac{\partial}{\partial y}\Bigr),
  \qquad
  \frac{\partial}{\partial\bar z}=\frac12\Bigl(\frac{\partial}{\partial x}+i\frac{\partial}{\partial y}\Bigr)
\]
of Section~2 of the paper, applied in the distributional sense.  It is the
reading the paper's text forces where those operators are first used.  A
mapping is called $K$-quasiregular at paper.tex:129--133 if
$f\in W^{1,2}_{\mathrm{loc}}(\Omega)$ \emph{and} the inequality (11) between
$\partial f/\partial\bar z$ and $\partial f/\partial z$ holds almost
everywhere; at that point in the paper $f$ carries no hypothesis beyond
Sobolev membership, so the derivatives appearing in (11) --- and in (12),
which is about the same $f$ one sentence later --- are necessarily the
distributional ones whose existence that membership asserts.  The same is true
of the system (18), whose solutions the paper assumes only to lie in
$W^{1,2}$.

Clause (a) is part of the data for two reasons.  First, it is what makes
clause (b) say what it should: the integration-by-parts identity of
\noderef{IsWeakPartialOn} is an identity between integrals, and a divergent
integral carries no information, so without local integrability that identity
could be met by functions which are not derivatives of $f$ in any sense.
Second, (a) and (b) are exactly the hypotheses under which the weak partial
derivatives of an almost everywhere differentiable function are known to agree
almost everywhere with its classical ones.  Consequently, when $\Omega$ is
open and $f$ is in addition differentiable at almost every point of $\Omega$,
one has $f_x=\partial_0f$ and $f_y=\partial_1f$ almost everywhere, hence
$a=\partial f/\partial z$ and $b=\partial f/\partial\bar z$ almost everywhere
with the right-hand sides read through the real Fr\'echet derivative of $f$.
So for an almost everywhere differentiable function the weak and the pointwise
readings of the two operators agree almost everywhere, which is what lets the
pointwise computations of Sections~4--6 proceed.  For a function known only up
to a null set, by contrast, only the weak reading has content: the Fr\'echet
derivative is taken to be $0$ wherever a function fails to be differentiable,
so both pointwise operators vanish identically on a nowhere differentiable
representative of a class, and any inequality or identity between them is
satisfied vacuously there.

Finally, $a$ and $b$ are determined by $f$ and $\Omega$ up to null sets when
$\Omega$ is open: if $f_x$ and $f_x'$ are both distributional
$\partial_0$-derivatives of $f$ on $\Omega$ lying in
$L^1_{\mathrm{loc}}(\Omega)$, then $\int(f_x-f_x')\phi\,d\sigma=0$ for every
test function $\phi$ for $\Omega$, whence $f_x=f_x'$ almost everywhere on
$\Omega$, and likewise for $f_y$.  So the choice of the family in (b) does not
matter, and a statement of the form ``there are weak Wirtinger derivatives
$a,b$ of $f$ on $\Omega$ such that $P(a,b)$ holds almost everywhere'' is
unambiguous for every property $P$ which is insensitive to changes of its
arguments on null sets.
\end{definition}
